% Same-history physical stress transfer and assembly of trace chart data.
\subsection{The common-trajectory physical stress transfer}
\label{sec:ex-lattice-trace-scribe-producer}

This section proves the physical stress and flux passage used by the
selected trace chart family from the explicit response, Ward comparison, fourth-order,
and weighted trace receipts displayed below. Their same-history producers are identified and applied below; the
arguments also show how the resulting estimates combine. All limits below use one selected history, one
ultrafilter, one hierarchy, and one scaling. In particular, the physical
stress, the correction tensors, and the comparison observables are never
reselected on separate subsequences. The quantitative estimates entering the
construction include the paid response and comparison estimates from
Sections~\ref{sec:ex-lattice-trace-weighted-estimates},
\ref{sec:ex-lattice-trace-crown-buffer}, and
\ref{sec:ex-lattice-trace-eve-producer}. The additional Ward comparison and branch
receipts are stated explicitly below. Given these inputs, this section proves
how the estimates produce the conserved tensor sequence and transfer its
physical normalization to the selected table.

Write $\kappa=s_G>0$ for the slope of the actual clock. The chosen
charged response uses this same constant: its target payment is
$-12s_G$, and division by the flux normalization factor $12$ gives
$-s_G$. Thus the principal Ward limit below and the clock slope have
exactly the same $\kappa$. The final logarithmic trace coefficient is
independent of its numerical value, as shown below; no identification
with a separately normalized one-loop constant is needed here. Let $1$ denote the unit
field of the reference grade-four table, and let $\mathcal O$ be the
normalized curvature-square field of the reference table: the image in
that table of one quarter of the sum, over all ordered pairs $(a,b)$ of
coordinate directions, of the labels representing the color inner
products $\langle F_{ab},F_{ab}\rangle$ of the curvature with itself. It is the field denoted
$\mathcal O$ in Subsection~\ref{sec:ex-trace-branch-moments} and $O$ in
Subsection~\ref{sec:trace-asymptotic}. It has expectation zero in the
reference table, it is nonzero, and it is not a multiple of the unit
field. The physical coefficient is $d_2=-\kappa/4$.

\subsubsection{The two-direction conserved correction}

Let $E$ be the finite-dimensional real space of centered, signed-equivariant,
spin-symmetric grade-four tensors for the selected table (before imposing
the divergence equation), let $V_{\mathrm{int}}$ be the real space of
grade-five vector arrays of the selected table in which divergences lie,
let $V$ be the finite-dimensional response space, which is the grade-five
vector table of the reference table, and let $\mathcal B$ be the
grade-five profile space. Let
$D:E\to V_{\mathrm{int}}$ be the divergence map. Its codomain belongs to
the selected table, whereas the comparison and response maps below take
values in the reference space $V$. The conserved-line construction in
Subsection~\ref{sec:actual-stress-normalization}, applied to the completed
maps of Theorem~\ref{thm:ex-lattice-selected-trace-gallery}, gives
$\dim\ker D=1$ by \eqref{eq:stress-physical-line}. Its hypotheses
are checked for these same maps in the normalized-energy application
below: coherent unit-preserving maps, uniform short-test convergence,
and the reference rank calculation. In particular, rank one does not
require the correction sequence about to be constructed. The
interacting divergence remains fixed; the comparison varies the optical map,
and its positive scale and invertible chart preserve the kernel.

The construction solves the correction equations at each
large index, rather than taking an unspecified kernel vector. Let
$\mathsf K_i:E\to V$ be the selected comparison map, let
$s_i:\mathcal B\to V$ be the paid response map, and let
$L_i,\Omega_i,\Delta_i$ be the
principal and two correction directions. The correction map is
\[
 \mathsf S_i(c_0,c_1)=c_0\Omega_i+c_1\Delta_i,
 \qquad \mathsf B_i=\mathsf K_i\mathsf S_i.
\]
Let $q_i>0$ be the selected arch scale at the summit level (the
quantity $\zeta(k_i)$ of the clock paragraph below; it is compared with
the branch unit's effective coupling in
Lemma~\ref{lem:ex-lattice-scribe-normalization-ratio}) and $\rho_i$ its
positive response-size bound; these are the scales in the paid bounds below.
The optical comparison identity, which is an input here, states that on
the selected tail the comparison map is the divergence followed by the
scaled inverse chart: for every $Q\in E$ and every axis $a$,
\begin{equation}
 (\mathsf K_iQ)_a=q_ir_i^5\,U_{i,5}^{-1}\bigl((DQ)_a\bigr),
 \label{eq:ex-lattice-scribe-optical-identity}
\end{equation}
where $r_i>0$ is the chart radius and $U_{i,5}$ is the grade-five chart
isomorphism from reference fields to fields of the selected table. Since
$q_ir_i^5>0$ and $U_{i,5}$ is invertible, $\mathsf K_iQ=0$ if and only if
$DQ=0$ on that tail.
For $a=0,1$, the corrected response estimates compare the corresponding
column with the common-trajectory linear response on a paid forecast
$w_i^a\in\mathcal B$:
\[
 \|\mathsf B_i e_a-s_i(w_i^a)\|\longrightarrow0,
 \qquad w_i^a\longrightarrow v_a,
 \qquad s_i\longrightarrow\mathsf V.
\]
The profiles $v_0,v_1$ are the two free correction directions, and
$\mathsf W(c_0,c_1)=c_0v_0+c_1v_1$. Hence
\[
 \mathsf B_i\longrightarrow
 \mathsf B_\infty=\mathsf V\mathsf W,
 \qquad
 \mathsf V\mathsf W(c_0,c_1)=c_0\mathsf V(v_0)+c_1\mathsf V(v_1),
\]
where $\mathsf V$ is the free vector-response map and $\mathsf W$ is the
two-column reference correction. Indeed, for each $a$, the displayed
response estimate and continuity of the limiting response give
$\mathsf B_i e_a\to\mathsf V(v_a)$. Since the domain has basis
$e_0,e_1$, convergence of both columns implies operator-norm convergence.
The map $\mathsf B_\infty$ is an isomorphism. To see this, use the real
coordinates $w=(w_0,w_1,w_2)\in\mathbb R^3$ that parametrize
$\mathcal B$ bijectively and in which the branch profiles are defined.
In these coordinates
\[
 v_0=(0,0,1),\qquad v_1=(0,\tfrac14,0),
\]
and two facts from the reference computation of
Subsection~\ref{sec:ex-lattice-trace-free-rank} are used as inputs:
$\mathsf V$ maps $\mathcal B$ onto $V$, and its kernel is the line
spanned by the equation-of-motion direction, whose coordinates are
$\epsilon=(1,-\tfrac14,1)$. If $\mathsf V\mathsf W(c_0,c_1)=0$, then
$c_0v_0+c_1v_1=\gamma\epsilon$ for some real $\gamma$; the coordinate
$w_0$ gives $\gamma=0$, then $w_2$ gives $c_0=0$ and $w_1$ gives
$c_1=0$. So $\mathsf V\mathsf W$ is injective. It is onto: an arbitrary
element of $V$ is $\mathsf V(w)$ for some $w$, and
\[
 w-w_0\epsilon=(w_2-w_0)v_0+(4w_1+w_0)v_1,
\]
so $\mathsf V(w)=\mathsf V\mathsf W(w_2-w_0,\,4w_1+w_0)$ because
$\mathsf V(\epsilon)=0$. In finite dimension, if
$\|\mathsf B_i-\mathsf B_\infty\|<1/(2\|\mathsf B_\infty^{-1}\|)$,
then for every $c$,
\[
 \|\mathsf B_i c\|\ge
 \|\mathsf B_\infty^{-1}\|^{-1}\|c\|-\|
   (\mathsf B_i-\mathsf B_\infty)c\|
 \ge \frac{\|c\|}{2\|\mathsf B_\infty^{-1}\|}.
\]
Thus $\mathsf B_i$ is injective and hence bijective, and its inverse norms
are uniformly bounded on a tail. Let $X_i=q_i\mathsf K_i(L_i)$ be the
uncorrected forcing. The paid response-size estimate gives
$\|X_i\|\le C\rho_i$, where $\rho_i$ is the positive response-size
bound. Define
\[
 c_i=-\mathsf B_i^{-1}X_i,\qquad
 T_i=q_iL_i+\mathsf S_i(c_i).
\]
Then $\mathsf K_i(T_i)=0$, and the optical comparison identity
\eqref{eq:ex-lattice-scribe-optical-identity} gives
$DT_i=0$: on the selected tail, $\mathsf K_i(Q)=0$ is equivalent to
$DQ=0$ for every $Q\in E$. The inverse bound gives
$\|c_i\|\le C'\rho_i$.

The scalar calibration uses the same paid functional $\mu_i\in V^*$ that controls
the response bound. Let $\Lambda\in\mathcal B^*$ be the free Ward functional.
At chart level $k_i$, write $t_i$ for its positive Ward comparison
scale and $z_i=\widehat\beta_{k_i}$ for its clock weight. The bound comparison
gives $t_i\to\infty$, $\rho_i\le t_i\le z_i$. The Ward comparison inputs
include a constant $C_\mu$ with $\|\mu_i\|\le C_\mu$ eventually, and, for
every $\delta>0$, the eventual estimate
\[
 t_i\left|\mu_i(s_i y)-\Lambda(y)\right|
       \le\delta\|y\|\quad\text{for all }y\in\mathcal B.
\]
These Ward comparison inputs are the conclusion of
Proposition~\ref{prop:trace-prosper-sensor}. We use it under the
following identification of its notation with the present one. Its spaces $E$ and $H$ are $\mathcal B$ and $V$; its Ward comparison
maps $s_i$ are the present $s_i$, defined by the selected chart formula
in \eqref{eq:trace-sensor-selected-formulas} and not
merely maps with the
same limit; its free projection $f$ is $\mathsf V$; its Ward functional
$\rho$ is $\Lambda$; its postselected level $n_i$ is the present $k_i$,
the summit index of the frame level of the $i$th chart; and its weight
$\mathfrak a_i$ is $t_i$. In the notation of
Subsection~\ref{sec:ex-lattice-trace-weighted-estimates}, with $p_i=z_i$,
one has $t_i=A_{1/5}(i)$ and $\rho_i=A_{1/6}(i)$, and $q_i$ is the arch
scale $b_i$ at the level $k_i$. The required arch ratios are verified
below from the running estimate and the first-crossing schedule. The
additional schedule and column bounds in the Ward comparison proposition's hypotheses (S1) and (S6)
follow from \eqref{eq:trace-extra-stage-payment}, taking $\mathfrak C_j$ to
dominate the grade-five column parameters of the fixed stage profiles.
These are the same column budgets in the raw-step estimate, not
operator norms of a coupling-dependent response. The stage entry is
purchased before selecting the history and hence before its charts. For its
marked vector use the actual medium
$m_{k_i}=q_i^{-1}\mathsf q_{k_i}-\mathsf w_{k_i}$, with the chosen
and reference responses of
Subsection~\ref{sec:ex-lattice-trace-weighted-estimates}. That section's
weighted assembly gives $m_{k_i}\to x$ and
$t_i\Lambda(m_{k_i})\to0$; its chart-field decay gives
$z_i s_i(m_{k_i})\to0$. Since $t_i\le z_i$, the latter implies
$t_i s_i(m_{k_i})\to0$. These are exactly the three marked-vector
limits of the Ward comparison proposition, with the same normalized chart map.
The selected-lift formulas there use the same finite-volume inverse
column and the same fixed-level ultrafilter limit as this construction;
thus their identity $PY_i=X_i$ holds before and after projection,
not merely asymptotically. Apply its initial-level bound and transport
bound to this lift, taking the finite intersection of the levelwise
ultrafilter sets between the frame level and its summit. This is
legitimate at each fixed chart index; no tail uniform over all chart
indices is required. Its convergence-of-selected-lifts lemma therefore
applies to the same lift. The actual tested-bound and raw-step bounds (S4) and (S5) for
this lift are proved in the tested-bound and adjacent-column arguments
below. The identity and finite intersection only synchronize them. With
these quantitative bounds, the
Ward comparison proposition supplies the
bounded functionals $\mu_i$ and the estimate uniform over every
$y\in\mathcal B$ used above. We spell out how these conclusions
give the three limits needed for calibration. Let $Y_i$ be the principal
forecast. The paid response comparison, penumbra estimate, and principal
Ward limit are
\begin{equation}
 q_i\mathsf K_i(L_i)-s_i(Y_i)\longrightarrow0,\qquad
 \|Y_i\|\le C\rho_i\le Ct_i,\qquad
 \Lambda(Y_i)\longrightarrow-\kappa.
 \label{eq:ex-lattice-scribe-principal-response-input}
\end{equation}
The first convergence and the uniform bound on $\mu_i$ imply
$\mu_i(X_i)-\mu_i(s_i(Y_i))\to0$. For any $\delta>0$, the Ward comparison estimate
then gives eventually
\[
 |\mu_i(s_i(Y_i))-\Lambda(Y_i)|
       \le \delta\|Y_i\|/t_i\le C\delta.
\]
Taking the limit superior and then $\delta\downarrow0$ proves
$\mu_i(X_i)\to-\kappa$; no norm convergence of the principal forecast is
needed.

For a correction branch $P\in\{\Omega,\Delta\}$, let $U_i^P\in\mathcal B$
be its finite-volume profile and $v_P\in\mathcal B$ its reference profile.
The comparison and
Ward estimates needed are
\begin{equation}
 U_i^P\to v_P,\qquad
 z_i\|\mathsf K_i(P_i)-s_i(U_i^P)\|\to0,\qquad
 t_i|\Lambda(U_i^P)-\Lambda(v_P)|\to0,
 \label{eq:ex-lattice-scribe-correction-response-input}
\end{equation}
where $P_i$ denotes the corresponding correction tensor. Since $U_i^P$ is
bounded, the uniform Ward comparison estimate implies
$t_i|\mu_i(s_i(U_i^P))-\Lambda(U_i^P)|\to0$. Using $t_i\le z_i$ and the
uniform operator bound for $\mu_i$ in the middle estimate of
(\ref{eq:ex-lattice-scribe-correction-response-input}) gives
\[
 t_i\left|\mu_i(\mathsf K_i(P_i))-\Lambda(v_P)\right|\longrightarrow0.
\]
The reference Ward values are $\Lambda(v_\Omega)=0$ and
$\Lambda(v_\Delta)=-4$. Thus the principal and two correction limits are
now established from the explicitly stated paid inputs:
\begin{equation}
 \mu_i(X_i)\to-\kappa,\qquad
 t_i\mu_i(\mathsf K_i\Omega_i)\to0,\qquad
 t_i\bigl(\mu_i(\mathsf K_i\Delta_i)+4\bigr)\to0.
 \label{eq:ex-lattice-scribe-sensor-limits}
\end{equation}
Applying $\mu_i$ to the exact cancellation $X_i+\mathsf B_i c_i=0$ gives
\[
 \mu_i(X_i)+c_i^0\mu_i(\mathsf K_i\Omega_i)
       +c_i^1\mu_i(\mathsf K_i\Delta_i)=0.
\]
Set $\alpha_i=t_i\mu_i(\mathsf K_i\Omega_i)$ and
$\beta_i=t_i(\mu_i(\mathsf K_i\Delta_i)+4)$. Since
$\|c_i\|\le C\rho_i\le Ct_i$, both $c_i^0/t_i$ and $c_i^1/t_i$ are
bounded. The exact equation rewrites as
\[
 0=\mu_i(X_i)-4c_i^1+
       \frac{c_i^0}{t_i}\alpha_i+\frac{c_i^1}{t_i}\beta_i.
\]
The last two terms tend to zero by
(\ref{eq:ex-lattice-scribe-sensor-limits}), so
$\mu_i(X_i)\to-\kappa$ implies
\[
 c_i^1\longrightarrow-\frac{\kappa}{4}.
\]
This is the stable-inversion and coefficient-extraction argument of
Subsection~\ref{sec:ex-lattice-trace-producer}.

The correction is built from three paid tensor families, the same at
every index: a principal tensor $L_i$ and two correction tensors
$\Omega_i,\Delta_i$. These families and the coefficients $c_i^0,c_i^1$
below are the tensors of that selected construction. Their field and
trace properties follow from the common-trajectory branch application below. Let
$\mathcal P_i$ be the scaled pullback of a stress tensor, and put
$\mathcal Y_i=q_i^{-1}\mathcal P_i$. The scaled pullback is the
chart map multiplied by its positive crown normalization; by definition,
$\mathcal P_i=q_i\mathcal Y_i$ on every tensor. Here
$q_i=\zeta(k_i)$ is the arch scale at the selected level (see the clock
paragraph below), and the crown normalization is $q_i$ times the fourth
power of the chart radius. The defining-limit identification
\eqref{eq:scribe-same-arch} below identifies $q_i$ with the branch
unit's effective coupling at $k_i$; the subsequent ratio lemma also
records the transfer for asymptotically equal normalizations. The defining corrected tensor is
\begin{equation}
 T_i=q_iL_i+c_i^0\Omega_i+c_i^1\Delta_i.
 \label{eq:ex-lattice-scribe-corrected-tensor}
\end{equation}
For the three tensor sequences $Q_i\in\{L_i,\Omega_i,\Delta_i\}$, let
their reference branches be $\nu_L=0$, $\nu_\Omega=-$, and
$\nu_\Delta=+$. Let $W_i^Q(a,b)$ be the ordered grade-four label for the
two tensor indices $a,b$ obtained by the finite-volume response of the
corresponding sequence. These ordered
labels are the entries of a finite array; no spatial sum is hidden in this
notation. Let $\mathsf M_i$ be the finite-volume fourth-order coefficient
map, obtained by applying the centered chart pullback to a grade-four label.
It is a continuous linear map
$\mathsf M_i:\mathcal L_4\to\mathcal F_4$, not a scalar moment matrix.
Each $W_i^Q(a,b)$ belongs to $\mathcal L_4$. For every sufficiently large
$i$ (so the selected level exceeds the fixed reference cutoff), symmetrizing
the two ordered indices gives the chart formula
\begin{equation}
 [\mathcal P_i(Q_i)]_{ab}=
 \frac12\left(\mathsf M_i(W_i^Q(a,b))+\mathsf M_i(W_i^Q(b,a))\right).
 \label{eq:ex-lattice-scribe-phase-formula}
\end{equation}
The selected coefficient convergence gives $W_i^Q(a,b)\to
W_{\nu_Q}(a,b)$, and the fourth-order maps converge in operator norm,
$\mathsf M_i\to\mathsf M_\infty$; for the present normalization this
uses $\vartheta_i\to1$ of
Lemma~\ref{lem:ex-lattice-scribe-normalization-ratio}, since $\mathsf M_i$
is $\vartheta_i^{-1}$ times the branch unit's moment map. Thus, entry by
entry in the finite
stress array,
\[
 \mathcal P_i(Q_i)\longrightarrow\mathcal B_{\nu_Q},\qquad
 (\mathcal B_\nu)_{ab}:=
 \frac12\left(\mathsf M_\infty(W_\nu(a,b))+
                \mathsf M_\infty(W_\nu(b,a))\right).
\]
This follows directly from continuity of the finite bilinear expression in
(\ref{eq:ex-lattice-scribe-phase-formula}); no limit is interchanged with
an infinite sum. The explicit free tensor calculation identifies these
arrays with the reference material tensors. Write
\[
 M_0=\mathcal B_0,\qquad M_\Omega=\mathcal B_-,\qquad M_1=\mathcal B_+.
\]
The false-branch tensor need not vanish; it is its coefficient
$c_i^0/q_i\to0$ that removes it from the normalized field limit. The three
branch limits are on the same history.

The three branch limits, and the trace statements used below, are the
conclusions of Subsection~\ref{sec:ex-trace-branch-moments}, under the following
identification. The tensors are $L_i=t^0_{k_i}$, $\Omega_i=t^-_{k_i}$
and $\Delta_i=t^+_{k_i}$, at the same level $k_i$ as above. That
section uses the map $\mathsf P_i=\omega_iU_i^{-1}$ with
$\omega_i=\beta_Z(k_i)r_i^4$, where $\beta_Z$ is its effective coupling,
whereas the present crown normalization uses the arch scale $q_i$ in
place of $\beta_Z(k_i)$. Both normalizations use the same fixed-level
limit of the running coupling. Indeed the arch is constructed by taking,
for each level $j$, the limit of $\beta_{N,p,j}$ along the selected
ultrafilter; this defines $\zeta(j)$. The branch unit's $\beta_Z(j)$ is
the limit of this same real-valued function along this same ultrafilter.
Uniqueness of a limit in $\mathbb R$ therefore gives
\begin{equation}\label{eq:scribe-same-arch}
                  \beta_Z(j)=\zeta(j)\qquad\text{for every }j.
\end{equation}
No further extraction or adjustment of the coupling is involved. With
\[
 \vartheta_i=\frac{\beta_Z(k_i)}{q_i}>0,
\]
one has $\vartheta_i=1$ for the actual constructed arch. More generally,
the scaled pullback $\mathcal P_i$ is the entrywise application of
$\vartheta_i^{-1}\mathsf P_i$, and $\mathsf M_i$ is its centered moment
map, which is $\vartheta_i^{-1}$ times that of $\mathsf P_i$. The
ordered labels are $W_i^Q=W^{b}_{k_i}$ for the branch $b=\nu_Q$, and
$M_0=M(w^0)$, $M_\Omega=M(w^-)$, $M_1=M(w^+)$.
Proposition~\ref{prop:ex-branch-moments} gives the three field limits in
\eqref{eq:ex-lattice-scribe-field-receipts};
Proposition~\ref{prop:ex-branch-washout} with exponent $p=1$ gives the
weighted trace limits for the principal and the false branch that are
used below, and gives nothing for the true branch;
and the trace formula \eqref{eq:ex-branch-three-traces} gives
$\operatorname{tr}M_1=4\mathcal O$. Their recording, transport and chart
estimates concern exactly this supplied history: at fixed $i$, use the
same finite-volume tensors and the same ultrafilter before passing to
$i\to\infty$. The three branch choices and sixteen tensor entries form
a finite set, so their eventual identities and estimates hold on one
common tail. This finite intersection does not prove the actual branch
coefficient bias or local-current identities, nor the three recording
estimates in \eqref{eq:ex-branch-book-current},
\eqref{eq:ex-branch-book-difference}, and
\eqref{eq:ex-branch-book-tensor}. Those estimates are required at
these actual labels. No separate history is selected for a branch. The
finite-array argument written out above repeats the
proof of the first of them under the same hypotheses, together with
$\vartheta_i\to1$ of Lemma~\ref{lem:ex-lattice-scribe-normalization-ratio}.

Those conclusions are stated for $\mathsf P_i$. They transfer to
$\mathcal P_i=\vartheta_i^{-1}\mathsf P_i$ because of the following lemma.

\begin{lemma}[The two normalizations agree asymptotically]
\label{lem:ex-lattice-scribe-normalization-ratio}
Use the scale ratio $\beta_Z(j)/\widehat\beta_j\to1$ from Subsection~\ref{sec:ex-trace-branch-moments} (proved in its scale-ratio paragraph, after the display
\eqref{eq:ex-branch-lift-beta-ratio}, from its running estimate and
fixed-scale limit), and the first clock limit
$\zeta(j)/\widehat\beta_j\to1$ of
\eqref{eq:ex-lattice-scribe-clock-inputs}, proved below from
Lemma~\ref{lem:rank-trace-running}. Here $\widehat\beta_j$ is the
same clock coupling in both statements: the branch unit and this section
use the hatted coupling of the one selected clock. Then $\vartheta_i\to1$.
Consequently, for $Q_i\in\{L_i,\Omega_i,\Delta_i\}$ with reference
tensors $M_0,M_\Omega,M_1$,
\begin{gather*}
 \mathsf P_i(Q_i)\to M\ \Longrightarrow\ \mathcal P_i(Q_i)\to M,
 \qquad
 \mathsf P_i(\Delta_i)\to M_1\ \Longrightarrow\
 \operatorname{tr}\mathcal P_i(\Delta_i)\to4\mathcal O,\\
 z_i\operatorname{tr}\mathsf P_i(Q_i)\to0\ \Longrightarrow\
 z_i\operatorname{tr}\mathcal P_i(Q_i)\to0 .
\end{gather*}
\end{lemma}
\begin{proof}
The summit index $k_i$ is a level not below the frame level $j_i$, and
$j_i\to\infty$, so $k_i\to\infty$. Hence both ratio limits hold along
$j=k_i$. The quantities $\beta_Z(k_i)$, $q_i=\zeta(k_i)$ and
$\widehat\beta_{k_i}$ are positive, and
\[
 \vartheta_i=\frac{\beta_Z(k_i)/\widehat\beta_{k_i}}
               {\zeta(k_i)/\widehat\beta_{k_i}}
 \longrightarrow\frac11=1 .
\]
In particular $\vartheta_i^{-1}\to1$, and $\vartheta_i^{-1}\le2$ for all
large $i$. Entrywise, $\mathcal P_i(Q_i)=\vartheta_i^{-1}\mathsf P_i(Q_i)$
in the finite product of copies of $\mathcal F_4^0$; the product of a
convergent scalar sequence and a convergent vector sequence converges to
the product of the limits, which gives the first implication. The trace is
linear and continuous on that finite product, so
if $\mathsf P_i(\Delta_i)\to M_1$, then
$\operatorname{tr}\mathcal P_i(\Delta_i)=\vartheta_i^{-1}
\operatorname{tr}\mathsf P_i(\Delta_i)\to\operatorname{tr}M_1
=4\mathcal O$ by the trace formula of the branch unit. Finally
$\|z_i\operatorname{tr}\mathcal P_i(Q_i)\|
=\vartheta_i^{-1}\|z_i\operatorname{tr}\mathsf P_i(Q_i)\|
\le2\|z_i\operatorname{tr}\mathsf P_i(Q_i)\|$ for large $i$, which tends
to zero when the right side does.
\end{proof}

Since $z_i=\widehat\beta_{k_i}$, the washout limits of the branch unit
with $p=1$ are exactly the hypotheses
$z_i\operatorname{tr}\mathsf P_i(Q_i)\to0$ for $Q_i\in\{L_i,\Omega_i\}$.
By the lemma, the field limits and the exponent-one principal and false
trace limits used below hold for $\mathcal P_i$ on this same history.
For the actual arch this also follows directly from
\eqref{eq:scribe-same-arch}, since $\mathcal P_i=\mathsf P_i$. The exact
identity for $\mathcal Y_i(T_i)$ below uses only
$\mathcal P_i=q_i\mathcal Y_i$ and linearity, so it is unaffected.

These are the field receipts used below:
\begin{align}
 \mathcal P_i(L_i)&\longrightarrow M_0,&
 \mathcal P_i(\Omega_i)&\longrightarrow M_\Omega,&
 \mathcal P_i(\Delta_i)&\longrightarrow M_1,
 &\frac{\rho_i}{q_i}&\longrightarrow0,
 \label{eq:ex-lattice-scribe-field-receipts}
\end{align}
The displayed ratio is the selected-scale washout input: the response bound
is bounded by the one-sixth subpower clock bound, whose ratio to the effective
coupling tends to zero along this scale. Since $M_0$ is the prescribed
material tensor and $|c_i^0|+|c_i^1|\le C\rho_i$, it follows that
$c_i^a/q_i\to0$ for $a=0,1$. Linearity gives
the exact identity
\[
 \mathcal Y_i(T_i)=\mathcal P_i(L_i)
   +\frac{c_i^0}{q_i}\mathcal P_i(\Omega_i)
   +\frac{c_i^1}{q_i}\mathcal P_i(\Delta_i).
\]
The principal, false-correction, and true-correction rescaled tensors converge
to $M_0,M_\Omega,M_1$, respectively.
In particular the two correction phase tensors are bounded, and the exact
identity proves
\begin{equation}
 \mathcal Y_i(T_i)\longrightarrow M_0.
 \label{eq:ex-lattice-scribe-field-limit}
\end{equation}
For the scalar trace, the true correction profile is the true material
correction, whose diagonal trace is $4\mathcal O$ by the trace formula of
the branch-moment unit recalled above: the diagonal sum is
$4(0+4\cdot\frac14)\mathcal O=4\mathcal O$. We also use the scale bounds
$q_i\le2z_i$ and $\rho_i\le t_i\le z_i$ and the limits
\[
 z_i\operatorname{tr}\mathcal P_i(L_i)\to0,\quad
 z_i\operatorname{tr}\mathcal P_i(\Omega_i)\to0,\quad
 |c_i^0|\le C\rho_i,\quad
 \operatorname{tr}\mathcal P_i(\Delta_i)\to4\mathcal O,
 \quad c_i^1\to-\kappa/4.
\]
The first two limits are the exponent-one principal and false trace limits
of the branch-moment unit recalled above. The bound on $c_i^0$ is contained
in the coefficient bound $|c_i^0|+|c_i^1|\le C\rho_i$ used above. The
fourth limit follows from the field receipt
$\mathcal P_i(\Delta_i)\to M_1$ and $\operatorname{tr}M_1=4\mathcal O$, by
continuity of the trace. The coefficient limit $c_i^1\to-\kappa/4$ was
proved above, after (\ref{eq:ex-lattice-scribe-sensor-limits}), from the
exact cancellation $X_i+\mathsf B_ic_i=0$, the Ward comparison limits
(\ref{eq:ex-lattice-scribe-sensor-limits}), the coefficient bound
$\|c_i\|\le C\rho_i$ and the scale bound $\rho_i\le t_i$; it is therefore
conditional on the same inputs. These scale inequalities need no additional receipt. In the notation
of Subsection~\ref{sec:ex-lattice-trace-weighted-estimates},
$\rho_i\le\exp((\log\log z_i)^2/6)$ and
$t_i=\exp((\log\log z_i)^2/5)$, so $\rho_i\le t_i$ and
$t_i\to\infty$. Since $(\log\log z)^2=o(\log z)$,
$t_i/z_i\to0$. Finally the first clock input in
\eqref{eq:ex-lattice-scribe-clock-inputs}, applied at $k_i$, gives
$q_i/z_i\to1$, hence $q_i\le2z_i$ eventually and
$\rho_i/q_i\to0$. These consequences still depend on the stated
clock and response-size inputs.
To see exactly how the weighted estimates enter, use
$\mathcal P_i(T_i)=q_i\mathcal P_i(L_i)
+c_i^0\mathcal P_i(\Omega_i)
+c_i^1\mathcal P_i(\Delta_i)$. The principal contribution tends
to zero since $q_i\le2z_i$ and the exponent-one principal trace estimate
gives $z_i\operatorname{tr}\mathcal P_i(L_i)\to0$. For the false
correction, $|c_i^0|\le C\rho_i\le Cz_i$ and the exponent-one
weighted false-sector estimate gives
$z_i\operatorname{tr}\mathcal P_i(\Omega_i)\to0$.
Finally $c_i^1\to-\kappa/4$ and
$\operatorname{tr}\mathcal P_i(\Delta_i)\to4\mathcal O$.
Thus applying the trace to (\ref{eq:ex-lattice-scribe-corrected-tensor})
and using these three contributions gives
\begin{equation}
 \operatorname{tr}\mathcal P_i(T_i)\longrightarrow-\kappa\mathcal O.
 \label{eq:ex-lattice-scribe-flux-limit}
\end{equation}
Indeed, the principal term tends to zero, the first correction term is
bounded by $C z_i|\operatorname{tr}\mathcal P_i(\Omega_i)|$, and the last
term tends to $(-\kappa/4)(4\mathcal O)$. This establishes the physical
field and flux limits from the exact three-term correction, with no
interchange of limits or choice of history.

\subsubsection{Proportionality transfers the physical normalization}

We now show that the corrected physical tensor and the selected table
stress have the same limiting normalization. Let $A$ be the selected
algebra, and let $R$ be its centered table stress
$R=\operatorname{reCenter}_S(A.\operatorname{energy})\in E$, whose entries
are conserved and whose flux is the physical Ward flux. Let $g_i=T_i$ be
the explicit conserved correction sequence just constructed, viewed in the
same one-dimensional kernel of $D$. Its normalized field
and flux limits are (\ref{eq:ex-lattice-scribe-field-limit}) and
(\ref{eq:ex-lattice-scribe-flux-limit}). We prove the transfer, spelling
out the nonzero test and the scalar ratio.

The argument uses the following inputs, all for the same selected table,
reference table, derivative maps, chart and scaling as above; the
derivative maps of the two tables are assumed to obey the derivative and
orientation laws of their tables. None of these inputs
is a consequence of rank one or of the field and flux limits.
\begin{enumerate}
\item \emph{Physical stress laws.} The tensor $R$ lies in $E$, satisfies
 $DR=0$, and obeys the translation Ward identity of the selected table:
 for every separated configuration consisting of a distinguished point and
 spectator points, every tuple of spectator labels, every cutoff that
 separates the distinguished point from the spectator points, every axis
 $b$ and every field $B$, the shell measurement of $R$ with that cutoff
 and with $B$ at the distinguished point equals the correlation at that
 configuration with the derivative field $D_b^SB$ at the distinguished
 point and the same spectators.
\item \emph{Reference stress laws.} The material tensor $M_0$ is a
 centered, conserved stress tensor of the reference table and obeys the
 translation Ward identity of the reference table.
\item \emph{A detecting derivative.} There are an axis $a$, a reference
 observer $o$, and a reference grade-four field $v$ whose derivative in
 direction $a$ has nonzero value $y$ on the grade-five test of $o$.
\item \emph{Tested convergence and jet comparison.} The following four
 statements hold for the selected chart maps $U_{i,d}$ and radii $r_i$.
 \begin{enumerate}
 \item[(4a)] For every integrated reference test, that is, a compactly
  supported test function with fixed spectator labels and one
  distinguished slot, and every fixed reference field $g$ of the grade
  of that slot, the correspondingly rescaled tested correlation of the
  selected table with $U_{i,d}g$ in the distinguished slot converges to
  the reference tested correlation with $g$ in that slot. This holds
  also for the test obtained by differentiating the test function in
  the distinguished variable.
 \item[(4b)] There are finitely many integrated reference tests with a
  grade-five distinguished slot whose reference functionals have zero
  common kernel on the reference grade-five field space.
 \item[(4c)] In both tables a tested correlation with the derivative
  field $D_a g$ in the distinguished slot equals the derivative in
  direction $a$, in the distinguished variable, of the correlation with
  $g$ in that slot.
 \item[(4d)] Pointwise jet comparison for the observer $o$. For every
  fixed reference grade-five field $g$, the selected correlation at the
  configuration $r_ix$ with $U_{i,5}g$ at the distinguished point and
  the spectator labels $\operatorname{pins}_i\circ q$ converges to the
  observer's reference value $o.test_5(g)$. For every fixed reference
  stress tensor $H$, the rescaled shell measurement
  $\mathcal L_i(H)=\mathscr F_i(r_i^{-4}U_{i,4}H)$ defined below, in
  which the chart is applied entrywise, converges to the reference shell
  measurement $\operatorname{shellGrad}(H,\chi,a,x,q,v)$. Both are
  obtained below from the comparison-of-jets estimate for the finitely
  many grades and arities fixed by the observer, on a comparison window
  containing the support of the gradient of $\chi$ together with the
  spectator points. The deduction from that compact comparison is given
  below; the comparison must concern these same chart maps.
 \end{enumerate}
\end{enumerate}
The physical stress in item 1 is supplied by the following choice of the
energy in the selected algebra and the centering argument after it.
The reference laws and detector in items 2--3 are verified below for
the actual Wick comparator used by this chart family.
For the completed charts of the selected chart family, the comparison
statements in item 4 follow below from their uniform short-test estimate
and the derivative laws of the two tables.

\subparagraph{Identifying the actual reference fields and their normalization.}
We verify the reference identifications needed here before using the
Ward calculation. Keep the frame register of the supplied family fixed.
Let $c_d$ be its coordinate isomorphism from invariant jet polynomials
of grade at most $d$ to coordinate labels, let $\pi_d^0$ be projection
to reference fields, and let $\operatorname{obs}_d$ be passage to the
on-shell Wick polynomial. The actual reference table is defined by
\begin{equation}\label{eq:scribe-reference-dictionary}
 S_0\bigl(\pi_{d_1}^0c_{d_1}p_1,\ldots,
             \pi_{d_n}^0c_{d_n}p_n;x\bigr)
 =\operatorname{Wick}_x\bigl(\operatorname{obs}_{d_1}p_1,
                    \ldots,\operatorname{obs}_{d_n}p_n\bigr).
\end{equation}
Indeed its coordinate correlation is the Wick correlation composed in
each slot with $c_d^{-1}$; projection quotients precisely the labels
annihilated by all separated spectator correlations. Thus this identity
uses the reference table of the supplied family, not a replacement
table. Its Wick covariance is the curvature covariance
\eqref{eq:maxwell-ward-covariance}, obtained by taking the two curls
of the potential kernel $(4\pi^2|x|^2)^{-1}$; this normalization is part
of the defining free table. Write $\Phi_dp=\pi_d^0c_dp$. The on-shell polynomial address of
$\Phi_dp$ is $\operatorname{obs}_dp$. The free separation result of
Subsection~\ref{sec:ex-lattice-all-grade-free-detectors} makes this address an
isomorphism onto the invariant on-shell sector: its kernel is exactly
the correlation null space. In particular the reference is faithful.

Its derivative plan is obtained by conjugating polynomial translation
through this address. On a curvature jet, translation raises the jet
in the specified direction; on products it obeys the Leibniz rule.
The rotation action is obtained by conjugating rotation of the same
jets, including their tensor indices. Consequently, writing
$\delta_a$ and $\mathscr D_A$ for these polynomial derivations,
\begin{equation}\label{eq:scribe-reference-derivatives}
 D_a^0\Phi_dp=\Phi_{d+1}(\delta_ap),\qquad
 K_A^0\Phi_dp=\Phi_d(\mathscr D_Ap).
\end{equation}
These identities follow by applying the injective polynomial address
to either side. They are compatible with grade inclusions since those
inclusions preserve the addressed polynomial. They also prove the
positional derivative identity by differentiating the finite Wick
sum in \eqref{eq:scribe-reference-dictionary}. The reference data
installed before the comparison curve have exactly these derivative
and rotation maps; fixing the stress normalization preserves both maps
and the vacuum. Even if a derivative plan is specified instead by its
positional derivative law, it must be this one: subtracting the two
laws gives zero separated correlations for their difference on every
label, so faithfulness makes the difference zero. This verifies (R1)
and (R2) of Subsection~\ref{sec:free-maxwell-ward-flux} on these same data.

For the stress identification there is a sign to retain. The coordinate
pair $P_{ab,cd}$ in \eqref{eq:ex-branch-material} is the representative
of $\langle F_{ab},F_{cd}\rangle$, with no extra factor. The canonical
positive Maxwell tensor is therefore the image of
\eqref{eq:maxwell-wick-stress}; this verifies (R3). But the branch
coefficients $w^0=(-1,1/4,-1)$ in \eqref{eq:ex-branch-coefficients} give
\begin{equation}\label{eq:scribe-reference-stress-sign}
 (M_0)_{ab}=-\Phi_4\left(
   \sum_c\langle F_{ac},F_{bc}\rangle
      -\frac{\delta_{ab}}4\sum_{c,e}
                         \langle F_{ce},F_{ce}\rangle\right)=-J^0_{ab}.
\end{equation}
The products inside $\Phi_4$ are Wick ordered by the defining observable
map. Hence $M_0$ is centered and conserved. Applying the proved
polynomial Ward identity \eqref{eq:maxwell-polynomial-ward} and
\eqref{eq:scribe-reference-derivatives}, the outward collar flux of
$J^0$ is minus the derivative insertion, and that of $M_0=-J^0$ is
the derivative insertion with coefficient one. The same sign gives
coefficient one for rotations. This proves input~2; no choice depending
on the comparison sequence has entered it.

Finally the definition of the scalar label in the branch table gives
\begin{equation}\label{eq:scribe-reference-scalar-identification}
 \mathcal O=\Phi_4\left(\frac14\sum_{a,b}
                                  \langle F_{ab},F_{ab}\rangle\right).
\end{equation}
Here the sum is over ordered pairs, just as in
\eqref{eq:scribe-reference-stress-sign}. Applying
\eqref{eq:scribe-reference-dictionary} to two such labels, the explicit
Wick computation \eqref{eq:trace-free-scalar} proves
\eqref{eq:trace-reference-normalization} for this actual $\mathcal O$:
$S_0(\mathcal O,\mathcal O;x,y)=3\dim\mathfrak g/(\pi^4|x-y|^8)$.
Its expectation is zero by Wick ordering. For $\dim\mathfrak g>0$
this two-point function is nonzero, so $\mathcal O\ne0$; its zero
expectation also excludes a nonzero multiple of the unit. The
differentiated two-point test below now supplies input~3 without an
additional reference-normalization assumption.

\subparagraph{The tested-bound input for this lift.}
The Ward comparison's (S4) is supplied by the quantitative invoice of
Lemma~\ref{lem:trace-quantitative-invoice}, not merely by the convergence
of tested correlations. Indeed take its grade-five reading $T$ at the
same station, finite-volume cutoff and period as the selected lift.
Its calibrated map is $A=d_k^{-5}\operatorname{Col}_k$. Both maps send
the reference unit to the coefficient unit, and $A$ is injective by the
calibrated column isomorphism. The invoice gives, for every coefficient
vector $v$,
\[
 |Av-Tv|_1\le H^{-1/2}|Av|_1.
\]
Its threshold is uniform over the finite volumes and the chart offset;
it was paid in the same joint chart purchase. Enlarging the selected
tail gives the half-accuracy estimate as well. Thus the two estimates
in \eqref{eq:trace-sensor-tested-price} hold with the actual invoice bound
$\eta_i=H^{-1/2}$. In particular $T$ is invertible, and the selected
finite-volume lift is exactly
$T^{-1}\operatorname{Dil}_{\xi^{-1}}
 \operatorname{Pos}_{\sigma_i}(Uh)_\alpha$.
To check that the invoice remains valid after deleting the vacuum
coefficient, apply both bounds to
$v^\circ=v-(Tv)_0e_{\rm ref}$. Then
$Tv^\circ=\operatorname{core}Tv$ and
$\operatorname{core}Av^\circ=\operatorname{core}Av$, while
$|Av^\circ|_1\le2|Tv^\circ|_1$. Consequently
\[
 |\operatorname{core}Av-\operatorname{core}Tv|_1
 \le2\eta_i|\operatorname{core}Tv|_1.
\]
This is precisely the centered coefficient estimate needed for the
Ward comparison's initial-level bound, with the same inverse reading and lift.

\subparagraph{The adjacent-column receipt.}
Fix a level $m$ beyond the grade-five entrance, and retain the actual
calibrated columns at levels $m,m+1$ of one finite-volume hierarchy.
Write $L=L_m$, $B=C_m$, $U=C_{m+1}$, and let $k\ge1$ be the selected
local forecast bound and $r>0$ the fixed local remainder transport
bound at that transition. Put
\[
 z=|\operatorname{core}\widetilde z_{m+1}|_1,
 \qquad y=|\operatorname{core}\widetilde z_m|_1.
\]
Deleting the vacuum coefficient before applying the calibrated front
estimate gives $y\le UL^5z$. The block map obeys
$|\operatorname{Block}_m\operatorname{core}
\widetilde z_{m+1}|_1\le L^5z$.

Apply the actual insertion and local uniqueness comparison of
Subsection~\ref{sec:ex-lattice-trace-price-recurrence}, specifically
\eqref{eq:ex-lattice-price-one-step} and
\eqref{eq:ex-lattice-local-uniqueness}, to these adjacent calibrated
fields. Their format and remainder bounds are the calibrated field
certificates, and their exact density transfer is the defining recursion
of this same hierarchy. The joint estimate takes the following form. For every $\epsilon>0$, on a sufficiently
large period tail, uniformly over $N>m$ and over the terminal label,
\[
 |\widetilde z_m-\mathcal L_m\widetilde z_{m+1}|_1
 \le s_m(z+Uz)
   +\epsilon\{By+k|\operatorname{Block}_m
                 \operatorname{core}\widetilde z_{m+1}|_1+rUz\},
 \qquad s_m\le A_{-1/2}(m+1).
\]
Here $\mathcal L_m$ is the actual local forecast, not a freely chosen
comparison map. To see the period quantifier, compare the calibrated
coarse density with the forecast of the calibrated fine density.
Both are exact representations of the same transferred density. Cancel
the common nonzero vacuum density and combine the two absolutely
summable regular remainders, changing the sign of one family. The local
section of their coefficient difference is exactly this combined
remainder. Its nonwrapping terms have zero local projection. Subtract a finite
partial sum of those terms before applying a member of the finite
separating probe family. The corresponding finite probe is exactly
zero. The probe comparison estimate bounds the remaining error by
its fixed derivative bound times the sum of the omitted nonwrapping
weights and the weights of wrapping supports. If the period exceeds
$10(h(4B_m+2)+2B_m)$, every wrapping support anchored at the insertion
has enough weighted support mass to pay $e^{-a h}$, where $a>0$ is
the fixed support exponent. Send the finite partial sum through all
nonwrapping terms: its omitted weights tend to zero by absolute
summability. This finite-partial-sum argument avoids differentiating
an uncontrolled infinite series. Sum the finitely many probe bounds
and apply the fixed coefficient recovery matrix. The resulting bound
on the coefficient difference is
$C e^{-a h}$ times the sum of the two remainder bounds. Choose $h$ so that
$Ce^{-ah}\le\epsilon$, then require the period to contain that collar
and the fixed forecast support. The constants and the period threshold
are independent of $N$ and the terminal label. The forecast remainder
is bounded by $k|\operatorname{Block}_m\operatorname{core}
\widetilde z_{m+1}|_1+rUz$; its coefficient error is $s_m(z+Uz)$.
These two bounds give the displayed inequality by the triangle
inequality. The selected transition bound is
$s_m=w_m A_{-1/2}(m+1)$ with $0<w_m\le1$, which gives its stated bound.

The forecast map is exactly the primary map of
Subsection~\ref{sec:primary-calibration-producer}, not a matrix selected
from its bounds. To verify its normalization, let $R$ be that section's
unscaled coefficient recovery, put $f=\vartheta_{m+1}$,
$c=\vartheta_m$, and write $(g_i,\ell_i)$ for the two grades of its
fixed coordinate basis. The local map before block transport is
\[
 \mathscr B=\mathscr S_c R\mathscr D_L\mathscr S_{f^{-1}},
 \qquad \mathscr B_{bi}=c^{\ell_b}f^{-\ell_i}L^{g_i}R_{bi}.
\]
This is the same full-matrix identity as in the primary calibration;
Subsection~\ref{sec:heading-history-handoff} explicitly uses that same
$R$. Its proved identity \eqref{eq:heading-triangular-receipt}, with
$u=L^{-1}$, gives
\[
 \mathscr B_{bi}=(c/f)^{\ell_i}\delta_{bi}\qquad(g_b\le g_i).
\]
There is no error term in this identity: the polynomial transport of
a weight-$a$ potential is $u^a$ times that potential plus a term of
weight less than $a$; the local section's lower-weight cancellation
annihilates every chain-rule term using the latter. The section then
reads the original homogeneous coordinate exactly. This is precisely
the argument proving the cited triangular receipt on the actual kernel.

The same section proves that $R$, and hence $\mathscr B$, commutes with signed
coordinate motions: change variables in the absolutely convergent
kernel sum, use the local section's covariance, and take its finite
jet. This establishes the symmetry for this very matrix, rather than
requiring a separate symmetric approximation. Every nonunit invariant
coordinate through grade five has factor degree two and position
degree four or five: each curvature factor costs at least two units
of position degree, and a linear invariant functional on a simple
Lie algebra vanishes. The only possible off-diagonal matrix entry left
by triangularity therefore raises position degree four to five.
Total inversion acts on degree $g$ by $(-1)^g$, so commutation forces
that entry to equal its negative and hence to vanish. The unit is
fixed exactly by the local map. Thus on the entire grade-five frame
$\mathscr B=\operatorname{Dil}_{c/f}$.

The full raw forecast is this same $\mathscr B$ composed with
$\operatorname{Block}_m$, whose coordinate weights are $L^{5-g_i}$.
Consequently
$\mathcal L_m=\operatorname{Dil}_{\vartheta_m/\vartheta_{m+1}}
\operatorname{Block}_m$. Its unit weight is $L^5$, consistently with
$\operatorname{Col}_m(e_{\rm ref})=a_m^5e_0$ and
$a_m=L a_{m+1}$; only the pre-block local map fixes the unit.
The selected threshold comparison gives
$\vartheta_{m+1}/\vartheta_m\le2$.
Set $g_m=2(1+U)A_{-1/2}(m+1)$ and
$P=BUL^5+kL^5+rU>0$. Choosing $\epsilon=g_m/(2P)$ in the preceding
period estimate bounds its second term by $(g_m/2)z$; its first term
is also at most $(g_m/2)z$. This proves exactly
\eqref{eq:trace-sensor-raw-step}, with its original calibrated columns.
The period tail is chosen separately for each fixed $m$, as required
in (S5). Intersecting finitely many such tails between the chart level
and its summit gives the single ultrafilter-large set used in the
Ward comparison's transport proof.

\subparagraph{The actual branch recording family.}
The recording estimates in the branch unit are applications of one
finite-volume purchase to a finite family of local observables. We
specify that family to keep its coefficient and rooted bounds attached
to the same functions. The principal tensor is the localized stress
field, the false correction is its off-diagonal curvature-product
array, and the true correction is the diagonal quarter-square array.
Include every tensor entry, its directional forward differences and
divergence, and the trace and antisymmetry currents of the localized
stress together with their forward differences. Tensor entries have
grade four, their differences and divergence grade five, the currents
grade three and their differences grade four. There are finitely many
branches, axes and pairs of tensor indices. Translations are applied to
the already chosen local functions and their labels; they do not
require a new purchase for each test point.

The local analytic construction and healthy derivative calculation of
Subsubsection~\ref{subsec:trace-charge-healthy-prices} apply to precisely these
entries, currents and differences. Subsubsection~\ref{subsec:trace-reference-atomic-certificate} supplies the
regular atomic certificate uniformly over this finite family in grades
at most five; its support-tail estimate applies to the original full
observable, including its healthy mask. Apply Subsection~\ref{sec:ex-lattice-trace-coffer-supplier} using the same
selected coffer for every member, and choose its coefficient label
once for each scale and finite-volume index, together with the full
purchase certificate. Its coefficient bias is
$\vartheta_k^2 A_{-1/4}(k)$. For the rooted estimate substitute
$a=\widehat\beta_k^{-1}$, $B=\widehat\beta_k^3$,
$\zeta_j=\varphi_{d,j}$, $\|\operatorname{invDil}_{j,d}\|=a_j^{-d}$
and $\operatorname{gain}(k,j)=e^{(s_k-s_j)/8}$ into
\eqref{eq:ex-trace-coffer-weight} and
\eqref{eq:ex-trace-coffer-root}. This gives exactly
\[
 \frac{2\operatorname{stake}_{h,j}a_j^{-d}}{\varphi_{d,j}}
       \widehat\beta_k^4e^{-(s_k-s_j)/8},
\]
which is the actual-minus-local bound in \eqref{eq:ex-branch-native-root-price}. The reference-frame column
bound gives a fixed-scale compact ball for the chosen labels. Hence
one ultrafilter limit $x_k$ exists for each member, before any
spectator test is chosen. This is the actual recording selection,
not a selection of a different label for each test.

Here is how the rooted receipt yields all three recorded inequalities.
Fix a spectator test and choose a native scale small enough that its
fixed points and a fixed neighborhood of the moving insertion are
separated. Its centered spectator bounds are bounded on that
neighborhood. Multiplying the actual-minus-local rooted bound by the
fixed spectator bound gives a bound uniform over the allowed integer
translations. At every fixed scale the point-limit identity for a
moving chosen label follows from its finite basis expansion: its
coefficients converge and the correlations for the fixed basis vectors
converge along the same ultrafilter. For each fixed $k$ and allowed displacement, choose the
ultrafilter-large set so that this convergence error is at most $a_k$.
The allowed integer displacements at that fixed level form a finite set,
so their intersection is still large. Adding $a_k$ to the rooted bound
gives the recorded current and difference inequalities; the tensor inequality
is the same argument at zero displacement. Their errors are fast:
with the native and comparison scales fixed, the rooted bound has the
form $C\widehat\beta_k^4e^{-s_k/8}$, so it is at most
$e^{-s_k/16}$ on a tail, since
$\log\widehat\beta_k=o(s_k)$. Multiplication by any fixed spectator
constant preserves such an exponential bound after a further tail.
Finite maxima over the family give one common tail. This proves the
actual instances of the three inequalities labelled
\eqref{eq:ex-branch-book-current},
\eqref{eq:ex-branch-book-difference}, and
\eqref{eq:ex-branch-book-tensor} in Subsection~\ref{sec:ex-trace-branch-moments}.

The current observable is normalized by $a_k^{-3}$ and its raw
first difference by $a_k^{-4}$. Hence the physical difference is
$a_k^{-1}$ times the difference of normalized current observables,
which is the factor in the external branch identity. The atomic
certificate applies to the raw local difference before this grade
normalization; no exponentially enlarged raw certificate is assumed.
The finite-volume identities are preserved in this application. The
trace of the principal tensor is the sum of the designated current
differences, and the false correction has exactly zero diagonal sum.
The trace and antisymmetry current identities and vanishing of their
projected forward differences are the local identities proved in
Subsection~\ref{sec:ex-lattice-trace-local-charge}; the correction entries
are the actual curvature products used in its local calibration.
Taking the finite sums before the ultrafilter limit therefore gives
exactly the current/difference relation and the principal and false
trace relations used by the branch washout proof. Its grade-three
unit argument, finite-difference estimate and moving finite detector
then apply to these recorded labels. Consequently the three branch
field limits and the principal and false weighted trace washout are
supplied for this very history and these very chart maps.

\subparagraph{Installing the normalized energy in the selected algebra.}
Keep the hierarchy, positional derivatives and orthogonal field action
constructed in Sections~\ref{sec:actual-stress-normalization},
\ref{sec:physical-ward-propagation} and
\ref{sec:physical-angular-actions}. Here is the same-table application.
Use the selected table $S$ as the physical hierarchy $S^1$ of the first
section, and use the actual comparator identified above as $S^0$.
The completed maps are those of the selected trace chart family; their
uniform test bound and derivative transport are verified below for
these same maps. The reference divergence is onto with one-dimensional
kernel by Subsection~\ref{sec:ex-lattice-trace-free-rank}. Thus
\eqref{eq:stress-physical-line} produces the conserved line in this
very selected table. Its choice of a nonzero tensor and sign uses a
subsequence only to establish existence of a fixed tensor $J$; it
does not change the table or its finite-volume extraction.

The annular-field construction of
Proposition~\ref{prop:ex-lattice-annular-field} applied to this $J$
gives its translation and angular response maps. The physical packet
bounds and their finite-row limits in
Subsection~\ref{sec:physical-ward-propagation} concern precisely these
response maps, so Corollary~\ref{cor:selected-ward-package} applies.
Its two free tests are supplied for the actual comparator by
Theorems~\ref{thm:free-quiet-classification},
\ref{thm:free-collinear-scenes} and
\ref{thm:free-internal-erasure}. No change of reference derivatives is
needed, by \eqref{eq:scribe-reference-derivatives}. The corollary gives
the global Ward laws with one nonzero scalar on the same $S$; the
inclusion argument makes that scalar common to all grades. Finally
Subsection~\ref{sec:physical-angular-actions} integrates those very
angular maps and makes them compatible with the selected derivatives
and inclusions. For this application, use the shell-node receipt at each opened pair
required by the finite watches. The actual schedule, lightness,
finite-stencil chamber placement and birth quota are verified in
\hyperref[par:shell-actual-placement]{the selected-family placement
argument}. The finite-mask identification is
\hyperref[par:shell-actual-mask]{the actual fine finite-mask payment},
and the coarse and fine transfer duties follow from
\eqref{eq:trace-extra-stage-payment} and its ensuing estimates.
These use the same selected service and interfaces as the two-channel
construction. The tab bounds follow from the adjacent clock ratio;
centered spectator bounds come from the one-field centered service.
For every fixed finite family of watches, the slotwise estimate
\eqref{eq:scribe-linear-derivative-price} supplies the required
separated-point derivatives with the enlarged spectator allowance.
It precedes the Ward construction and therefore introduces no use of
rotation covariance in proving that covariance.

The finite-row comparisons then use exactly this annular response,
the fixed normalization \eqref{eq:stress-fixed-normalization} and the
same-curve compact comparison. Their translation and rotation readout
limits are Proposition~\ref{prop:ward-readout-limits}; its reference
flux identity is the one identified above. The packet estimates are
those of Subsection~\ref{sec:physical-ward-packet-prices}, with the
stress-cone estimate of Subsection~\ref{sec:physical-ward-stress-cone}.
Their polynomial derivative input follows from the same pre-Ward
slotwise bound. Thus the shell, finite-row and packet arguments use
one selected hierarchy and its already fixed tensor throughout.
Let $J$ be the single conserved tensor in that construction and let
$c\ne0$ be its common charge scalar: its translation response is
$c\partial$ and its rotation response is $cD$, where $D$ is the Lie
algebra action of the constructed orthogonal representation $R_d$.
The energy of the selected algebra is chosen to be
\begin{equation}\label{eq:scribe-selected-energy}
 E_{ab}=\widehat J_{ab},\qquad
 \widehat J=\int_{O(4)}\mathcal R_K(c^{-1}J)\,dK,\qquad
 (\mathcal R_KQ)_{ab}=\sum_{i,j}K_{ai}K_{bj}R_4(K)Q_{ij}.
\end{equation}
The measure is probability Haar measure. This choice is made once,
together with the derivatives and the action, before choosing a trace
chart or a null sequence. It specifies the energy component of that
selected algebra; it is not an identification of an arbitrarily
preassigned energy with a generator of the conserved line.

Here is the same-data verification. The common scalar is independent of
grade by the inclusion argument of
Subsection~\ref{sec:physical-ward-propagation}. Multiplying the stress by
$c^{-1}$ multiplies each collar flux by $c^{-1}$, so its charges become
exactly $\partial$ and $D$. The averaging argument in the subsection
``Normalizing and averaging the stress'' of
Subsection~\ref{sec:physical-angular-actions} applies to these very maps.
It proves that the tensor in \eqref{eq:scribe-selected-energy} is real,
symmetric, tensor covariant and conserved, and that both Ward identities
retain coefficient one. Explicitly, for each $K$ the change of variables
$y=Kz$ carries its collar integral to that of $c^{-1}J$ with cutoff
$\chi\circ K$, transformed spectators and vector field $K^{-1}W(Kz)$.
Covariance of the same derivative and rotation action returns the
original Ward right-hand side. Each collar flux is a continuous linear
functional on the finite-dimensional space of stress arrays; it commutes
with the Haar integral. The right-hand side is independent of $K$, so
probability normalization leaves it unchanged. Thus the energy component
$E$ has the raw coefficient-one laws needed below.

No scale-dependent rescaling or second tensor is introduced in this
assembly. In the remainder of this section, the raw energy is exactly
\eqref{eq:scribe-selected-energy}, and $R$ is its centered version.
With those producer instances, the energy is installed on the selected
table with the raw Ward laws just obtained. No identification with an independently chosen
stress tensor is needed. In particular the auxiliary comparison
subsequence used to construct $J$ does not restrict the later trace
scalings: once constructed, $J$, its scalar and all its Ward identities
are fixed statements on $S$.

\subparagraph{Centering preserves the physical Ward normalization.}
There is no scalar-normalization loss in passing from the raw energy to
$R$. More precisely, suppose the raw energy $E_{ab}=A.\operatorname{energy}_{ab}$
is real, symmetric and tensor covariant, is conserved, and satisfies the
coefficient-one translation Ward identity
of input~1 on this same table, whose derivative annihilates the unit and
whose unit insertion removes that slot. Write
$c_{ab}=\langle E_{ab}\rangle_S$, so $R_{ab}=E_{ab}-c_{ab}\mathbf1$.
Then $\sum_a D_a^S R_{ab}=\sum_a D_a^S E_{ab}=0$.
For fixed spectator fields, the contribution of the subtracted tensor
to a translation shell is a sum of terms
\[
 -c_{ab}\,S(\text{spectator fields})
          \int_{\mathbb R^4}\partial_a\chi(z)\,dz=0.
\]
The last equality follows by integrating a derivative of a compactly
supported smooth function, one coordinate at a time. Thus the shell
identity for $R$ has exactly the same coefficient one and the same
derivative insertion as the identity for $E$. Centering is immediate;
symmetry, reality and tensor covariance are preserved because expectation
and the unit respect those structures. This proves the stress laws in
input~1 from the raw energy laws of the preceding construction, under
the shell and packet hypotheses specified there.
The scalar comparison with the explicit corrected tensor sequence
remains necessary and is proved below on the original trace scaling.

The detecting derivative in item 3 can also be fixed explicitly.
Take $v=\mathcal O$, the grade-four curvature-square field defined at
the start of this section, and choose any coordinate axis $a$. Let $o$
be the reference observer which tests a grade-five field $g$ by the
two-point correlation $S_0(g,\mathcal O;0,e_a)$. By the reference-field normalization
\eqref{eq:trace-reference-normalization}, the Wick calculation gives, with
$N=\dim\mathfrak g>0$,
\[
 S_0(\mathcal O,\mathcal O;x,e_a)
       =\frac{3N}{\pi^4|x-e_a|^8},\qquad
 o.test_5(D_a\mathcal O)
       =\left.\partial_{x_a}\frac{3N}{\pi^4|x-e_a|^8}
         \right|_{x=0}=\frac{24N}{\pi^4}\ne0.
\]
The derivative identity is the reference derivative law; both points
are separated, so ordinary differentiation is legitimate. This fixes a
single observer and a single nonzero value before any chart is selected.
It uses the actual reference Wick covariance and the identification
\eqref{eq:scribe-reference-scalar-identification} proved above. Thus the
detector is fixed without using the interacting trace asymptotic or
retaining an additional reference-identification premise.

\subparagraph{From chart comparison to the tested and shell limits.}
The tested approach property of
Subsection~\ref{sec:ex-lattice-firesky-gallery}, on the same scaling,
supplies (4a): a derivative of a fixed compactly supported smooth test
is another such test, with support in the same separated window.
That section's finite separating-test argument, applied in grade five,
supplies (4b). More explicitly, a nonzero field has a nonzero separated
reference evaluation. Continuity and a nonnegative smooth bump around
that evaluation give a nonzero integrated functional. These functionals
separate the grade-five space; finitely many suffice, since adjoining
a functional nonzero on the current common kernel strictly decreases
its dimension until it is zero. The derivative laws (4c) remain the
laws of the selected and reference tables.

We now check the stronger compact comparison needed in (4d). Use exactly
the completed corrected maps of
Theorem~\ref{thm:ex-lattice-selected-trace-gallery} and its
\hyperref[par:ex-trace-completed-uniform-tests]{completion paragraph}:
the index $k$ there is the present $i$, the selected table is $S$, the
reference table is $S_0$, and the radii and every grade map are the
present $r_i,U_{i,d}$. The raw lifts are those of these same supplied
charts and use the fixed extraction and ultrafilter. This is the
choice of chart family throughout the transfer, not a replacement by
a second curve. Any subsequent subsequence is applied to all these
data together.

For a fixed word $v$ and compact set $Q$ in its noncoincident
configuration space, the finite-slot expansion in that theorem gives
an error bound
\begin{equation}\label{eq:scribe-uniform-short-tests}
 \sup_{\substack{f\in C_c^\infty,\ \operatorname{supp}f\subset Q\\
                    \|f\|_\infty\le M}}
 \left|\int f(x)\left\{S(U_iv;r_ix)-S_0(v;x)\right\}\,dx\right|
       \le \frac{M K_v}{i+1}\longrightarrow0
\end{equation}
on a tail, where $K_v$ is a finite constant depending on the word.
Indeed the corrected discrepancy is a finite sum of raw discrepancies
at fixed words, each bounded by the short-test estimate
\eqref{eq:ex-lattice-trace-test-error}. Their number and label norms
do not depend on $f$ or $i$. The expanding cutoff window contains the
whole compact set $Q$ eventually. Completion changes none of the
finitely many grade maps on a further tail independent of $f$.
This proves the supremum estimate, rather than merely convergence for
each individual test. Along a further subsequence $i\mapsto n_i$,
the same argument retains the vanishing bound $MK_v/(n_i+1)$.

The compact-jet argument in the subsection ``Derivative transport and
convergence on compact sets'' of
Subsection~\ref{sec:actual-stress-normalization} now applies to these
very data. Its hypothesis \eqref{eq:stress-uniform-tests} is precisely
\eqref{eq:scribe-uniform-short-tests}. Its spaces $E_d,F_d$ are the
reference and selected grade spaces here; their finite dimension,
smooth multilinear correlations, separating tests, and derivative
identities are the same hierarchy and derivative laws already used
in (4a)--(4c). No stress-normalization conclusion of that section is
used in its compact-jet argument.

To make the application explicit, its transported derivative is
$r_iU_{i,d+1}^{-1}D_a^S U_{i,d}$. The finitely many separating tests
in grade $d+1$ give measurement maps $L_i\to L$, with $L$ injective.
Integration by parts and \eqref{eq:scribe-uniform-short-tests} imply
$L_i(r_iU_{i,d+1}^{-1}D_a^S U_{i,d})v\to LD_a^0v$.
The eventual lower bound for $L_i$ therefore gives convergence of
the transported derivative to $D_a^0$ in operator norm, as proved
there. Finite-basis multilinearity transfers the uniform test estimate
to converging moving labels. Position differentiation replaces one
label by its transported derivative, so this class of discrepancies
is closed under every fixed number of position derivatives. For a
compactly supported cutoff $\eta$ in configuration space and any such
discrepancy $h_i$, the tests $\eta e^{-2\pi\mathrm i\xi\cdot x}$
have one common supremum bound. The same is true with derivatives of
$\eta$ and descendants of $h_i$. The integration-by-parts calculation
in that section consequently bounds
$|\widehat{\eta h_i}(\xi)|$ by
$C\epsilon_i(1+|\xi|)^{-b}$, with $\epsilon_i\to0$ and any fixed
$b>4n$ at arity $n>0$. Fourier inversion yields uniform convergence
on a compact set where $\eta=1$; applying this to the finitely many
derivatives of a fixed order gives convergence of that jet. At arity
zero the assertion is the scalar short-test limit. This also explains
why no estimate at collision configurations is needed.

Thus, for each fixed finite word of reference
labels $h_1,\ldots,h_s$, with grades $d_1,\ldots,d_s$, and each compact
set $K$ of pairwise distinct configurations,
\begin{equation}\label{eq:scribe-compact-comparison}
 \sup_{z\in K}\left|
 S(U_{i,d_1}h_1,\ldots,U_{i,d_s}h_s;r_i z)
       -S_0(h_1,\ldots,h_s;z)\right|\longrightarrow0.
\end{equation}
The required finite grade bounds include the finitely many descendants
used for the desired jet order. They and the fixed arity are eventually
within the selected chart cutoff. Equation
\eqref{eq:scribe-compact-comparison} is therefore a proved consequence
of the uniform short-test estimate and the derivative laws, not an
independent compact-comparison premise.
Taking $K$ to be the observer's one configuration gives the first
assertion of (4d).

For the shell assertion, let $x$ denote the distinguished and spectator
points and choose the cutoff $\chi$ so that the compact set
$\operatorname{supp}\nabla\chi$ avoids all these points. The
configurations $(z,x)$ with $z\in\operatorname{supp}\nabla\chi$ then
form a compact subset of the noncoincident configuration space: all
pair distances have a positive minimum there. Thus one separated
window contains the entire integration set. Write $C_i^c(z;H)$ for
the selected correlation with the entry $U_{i,4}H_{ca}$ at $r_i z$,
the distinguished field $U_{i,4}v$ at its rescaled point, and the
chart-transformed spectators. Let $C_0^c(z;H)$ be its reference
counterpart. With the translation-shell convention, the flux is minus
the sum of these correlations against $\partial_c\chi$.
Changing variables $y=r_i z$
in its four-dimensional integral gives exactly
\begin{equation}\label{eq:scribe-shell-change}
 \mathcal L_i(H)
   =-\sum_{c=1}^4\int (\partial_c\chi)(z)C_i^c(z;H)\,dz.
\end{equation}
Indeed, the prefactor $r_i$ in $\mathscr F_i$, the stress multiplier
$r_i^{-4}$, the cutoff derivative $r_i^{-1}$ and the Jacobian $r_i^4$
multiply to one. The spectator maps are exactly the chart maps used
in \eqref{eq:scribe-compact-comparison}; no extra scaling of them is
introduced. That estimate on the compact shell therefore yields
\[
 |\mathcal L_i(H)-\operatorname{shellGrad}(H,\chi,a,x,q,v)|
 \le\sum_{c=1}^4\|\partial_c\chi\|_{L^1}
       \sup_{z\in\operatorname{supp}\nabla\chi}
          |C_i^c(z;H)-C_0^c(z;H)|\longrightarrow0.
\]
This proves the shell deduction in (4d). Expanding in a finite basis
of stress arrays also proves operator-norm convergence of these shell
functionals, which is the form used below for moving tensors.

By input~3 fix an axis $a$, an
observer $o$, and a reference grade-four field direction $v$ such that
$y=o.test_5(D_a v)\ne0$, where $D_a$ is the reference derivative in direction
$a$. Let $x$ be the separated configuration of this same observer,
consisting of the point of the distinguished insertion and the points of
its spectators, let $q$ be its fixed tuple of spectator labels, and let
$\chi$ be its cut, a cutoff separating the distinguished point from the
spectator points. The reference translation Ward identity of input~2, applied to
this cut, identifies this number with
$\operatorname{shellGrad}(M_0,\chi,a,x,q,v)$. We keep this exact observer and
cut for both stress sequences; no new test or window is selected in the
limit argument.

For a stress tensor $Q$, write the exact selected shell measurement as
\[
 \mathscr F_i(Q)=r_i\operatorname{shellGrad}_S\!\left(
 Q,\operatorname{stretch}(r_i,\chi),a,r_i x,
 \operatorname{pins}_i\circ q,U_i v\right),
\]
where $r_i$ is the chart radius, $q$ is the observer's tuple of spectator
labels, transported by the chart maps $\operatorname{pins}_i$,
and $U_i$ is the selected grade-four chart isomorphism. We claim
\begin{equation}
 f_i:=\mathscr F_i(R)\to y,
 \qquad \mathscr F_i(g_i)\to y.
 \label{eq:ex-lattice-scribe-common-test}
\end{equation}
The two limits have different sources, and both involve labels that
move with $i$, so that convergence for fixed labels cannot be applied
directly.

\emph{First limit.} Write $D_a^S$ for the derivative map of the selected
table and define the real-linear maps from reference grade-four fields to
reference grade-five fields
\[
 \mathsf T_i=r_i\,U_{i,5}^{-1}D_a^S\,U_{i,4}.
\]
We first show $\mathsf T_i\to D_a$ in operator norm. Let
$\lambda^1,\ldots,\lambda^m$ be the reference functionals of the tests
of (4b) and let $\lambda^1_i,\ldots,\lambda^m_i$ be the rescaled
functionals $g\mapsto$ (tested correlation of the selected table with
$U_{i,5}g$ in the distinguished slot). By (4a), $\lambda^r_i\to\lambda^r$
pointwise, hence in operator norm on the finite-dimensional reference
grade-five space. Fix a reference grade-four field $u$. By (4c) for the
selected table and integration by parts in the distinguished variable,
$\lambda^r_i(\mathsf T_iu)$ is minus the rescaled tested correlation of
$U_{i,4}u$ against the test with differentiated test function; the factor
$r_i$ in $\mathsf T_i$ is exactly the one produced by rescaling the
derivative. By (4a) for the differentiated test this converges to minus
the reference tested correlation of $u$ against the differentiated test,
which equals $\lambda^r(D_au)$ by (4c) for the reference table and
integration by parts. Thus $\Lambda_i(\mathsf T_iu)\to\Lambda(D_au)$,
where $\Lambda=(\lambda^r)_r$ and $\Lambda_i=(\lambda^r_i)_r$. Since
$\Lambda$ is injective on a finite-dimensional space there is $c>0$ with
$\|\Lambda w\|\ge c\|w\|$, and $\|\Lambda_iw\|\ge(c/2)\|w\|$ for large
$i$. Hence
\[
 \|\mathsf T_iu-D_au\|\le\frac2c\Bigl(
 \|\Lambda_i\mathsf T_iu-\Lambda D_au\|
 +\|(\Lambda-\Lambda_i)D_au\|\Bigr)\longrightarrow0,
\]
and pointwise convergence on the finite-dimensional grade-four space is
operator-norm convergence. Next let $A_i(g)$ denote the selected
correlation at $r_ix$ with $U_{i,5}g$ at the distinguished point and
spectator labels $\operatorname{pins}_i\circ q$. By the first part of
(4d), $A_i\to o.test_5$ pointwise, hence in operator norm. The physical
Ward identity of input~1, applied at the separated configuration $r_ix$
with the stretched cut, and linearity give
\[
 \mathscr F_i(R)=r_i\,(\text{correlation at }r_ix\text{ with }
 D_a^S(U_{i,4}v)\text{ at the distinguished point})
 =A_i(\mathsf T_iv).
\]
Since $A_i\to o.test_5$ in operator norm and
$\mathsf T_iv\to D_av$, we get
$\mathscr F_i(R)\to o.test_5(D_av)=y$.

\emph{Second limit.} No Ward identity for $g_i$ is used. For a reference
stress tensor $H$ put $\mathcal L_i(H)=\mathscr F_i(r_i^{-4}U_{i,4}H)$,
with the chart applied to each entry. This is a linear functional on the
finite-dimensional space of reference stress tensors, and by the
definition of $\mathcal Y_i$ and linearity of $\mathscr F_i$,
\[
 \mathscr F_i(g_i)=\mathcal L_i\bigl(\mathcal Y_i(g_i)\bigr)
\]
exactly. By the second part of (4d), $\mathcal L_i(H)\to\mathcal L_0(H)
:=\operatorname{shellGrad}(H,\chi,a,x,q,v)$ for every fixed $H$, hence
in operator norm. The field limit
(\ref{eq:ex-lattice-scribe-field-limit}) gives
$\mathcal Y_i(g_i)\to M_0$. Therefore
\[
 \mathscr F_i(g_i)\longrightarrow\mathcal L_0(M_0)
 =\operatorname{shellGrad}(M_0,\chi,a,x,q,v)=y,
\]
the last equality being the reference Ward identity used above.
The argument uses the same selected scale and the same observer test in both
limits. The physical flux excludes $R=0$: if $R$ vanished, linearity of
shell differentiation would force $f_i=0$ at every index, contrary to
$y\ne0$.

Since $\ker D$ is one-dimensional, $g_i=\lambda_iR$ for all sufficiently
large $i$. The second limit in (\ref{eq:ex-lattice-scribe-common-test})
therefore reads $\lambda_i f_i\to y$, while the first reads $f_i\to y$.
As $y\ne0$, choose a scalar coordinate functional $\ell$ with $\ell(y)\ne0$.
Then
\[
 \lambda_i=\frac{\ell(\lambda_i f_i)}{\ell(f_i)}\longrightarrow
 \frac{\ell(y)}{\ell(y)}=1.
\]
Thus $g_i=\lambda_iR$ and $\lambda_i\to1$. On a late tail
$\lambda_i\ne0$ and $\lambda_i^{-1}\to1$. Since $\mathcal Y_i$
and the weighted trace are linear,
\[
 \mathcal Y_i(R)=\lambda_i^{-1}\mathcal Y_i(g_i),
 \qquad
 \operatorname{tr}\mathcal P_i(R)=\lambda_i^{-1}
       \operatorname{tr}\mathcal P_i(g_i)
\]
on that tail. This proves the physical field and flux limits. In particular,
\begin{equation}
 \mathcal Y_i(R)\longrightarrow M_0,
 \qquad \operatorname{tr}\mathcal P_i(R)\longrightarrow-\kappa\mathcal O.
 \label{eq:ex-lattice-scribe-physical-transfer}
\end{equation}
This is the physical-array step used by the chart family. The dimension-one
kernel alone would not fix the scalar; the common nonzero shell test fixes it.

\subsubsection{Comparison tables and the complete trace chart data}

Let $T$ be any table on the same label book as the selected table $S$;
for the limit proved in this paragraph no hierarchy on $T$ and no
compatibility of $T$ with $S$ is required. Let
$\Theta_T$ be the trace of $A.\operatorname{energy}$ centered by its
one-point value in $T$. Its exact
centered decomposition is
\[
 \Theta_T=\operatorname{tr}(\operatorname{reCenter}_S(A.\operatorname{energy}))-c_T\,1_S,
 \qquad c_T\in\mathbb C,
\]
where $1_S$ is the grade-four unit field of $S$: both sides differ from
the uncentered trace by a multiple of $1_S$.
The physical-transfer limit (\ref{eq:ex-lattice-scribe-physical-transfer})
applies to the first term. The scalar-unit term vanishes after zooming:
the selected scale satisfies
\begin{equation}
 q_i r_i^4\longrightarrow0,
 \qquad q_i>0,
 \label{eq:ex-lattice-scribe-band-cost}
\end{equation}
as is proved from the clock in the next paragraph, and the selected chart
eventually identifies the reference unit with the unit of $S$, that is,
$U_i1=1_S$ for large $i$. Hence
$q_ir_i^4U_i^{-1}(c_T1_S)=c_Tq_ir_i^41\to0$. Therefore the zoomed
comparison observable has the same limit as the trace of the physical
stress,
\[
 q_ir_i^4U_i^{-1}\Theta_T\longrightarrow-\kappa\mathcal O=4d_2\mathcal O.
\]
The limit is a multiple of the curvature-square field $\mathcal O$, not of
the unit. This argument is
uniform in the algebraic identity for $T$; only the two scalar sequences
used below to represent its component in the distinguished two-dimensional
stress sector depend on $T$.

\emph{The clock.} Write $j_i$ for the frame level of the $i$th chart, so
that $k_i$ is its summit index, let $\mathsf a_j>0$ be the field scale at
level $j$ and $\mathsf b_i$ the zoom factor of the $i$th chart, and let
$\zeta(j)>0$ be the arch scale at level $j$; thus $q_i=\zeta(k_i)$. The
chart radius is $r_i=\mathsf a_{j_i}\mathsf b_i$, the running rate is
$\widehat\beta_j=\bar\beta+\kappa\log(\mathsf a_j^{-1})$ with a fixed
constant $\bar\beta$, and the selected run has $j_i\to\infty$,
$\widehat\beta_j\to\infty$ and $r_i\to0$. The first two clock limits
below follow from the actual arch construction. For each fixed $j$,
take the fixed-level ultrafilter limit in the running bound of
Lemma~\ref{lem:rank-trace-running}. Closedness of the real interval gives,
for every $\eta>0$ and every sufficiently large $j$,
\[
 |\zeta(j)-\widehat\beta_j|\le c_R+\eta s_j,
 \qquad \widehat\beta_j=\bar\beta+\kappa s_j,
 \qquad s_j=\log(\mathsf a_j^{-1}).
\]
The fixed constant $c_R$ does not depend on $j$. Division by the clock
and then $\eta\downarrow0$ proves $\zeta(j)/\widehat\beta_j\to1$.
For completeness, let $m+1=\operatorname{summit}(j)$ and
$x=\widehat\beta_j$. The selected schedule defines this summit as the
first crossing of the gap $(\log x)^{48}$. Thus
$0\le s_m-s_j\le(\log x)^{48}$ and
$\widehat\beta_m\le x+\kappa(\log x)^{48}\le2x$ eventually.
The actual schedule gives the weaker bound that is sufficient here:
by \eqref{eq:common-stage-entries},
\[
 0\le s_{m+1}-s_m=\log L_{J(m+1)}
       \le s_m/J(m+1).
\]
As $j\to\infty$, also $m\to\infty$ and $J(m+1)\to\infty$.
Since $s_m/x\le2/\kappa$ eventually, this last increment is $o(x)$.
The pre-crossing gap is $(\log x)^{48}=o(x)$, so
$\widehat\beta_{\operatorname{summit}(j)}/\widehat\beta_j\to1$.
In particular this clock-ratio proof does not need the stronger
block-size bound used in the external Ward comparison unit.
Since both levels tend to infinity,
\[
 \frac{\zeta(j)}{\zeta(\operatorname{summit}(j))}
 =\frac{\zeta(j)/\widehat\beta_j}
 {\bigl(\zeta(\operatorname{summit}(j))/
          \widehat\beta_{\operatorname{summit}(j)}\bigr)
  \bigl(\widehat\beta_{\operatorname{summit}(j)}/\widehat\beta_j\bigr)}
 \longrightarrow1.
\]
The third estimate is proved in
Lemma~\ref{lem:ex-lattice-trace-log-zoom} from the chart receipt
$1\le\mathsf b_i\le\operatorname{zoomCap}(j_i)$
of \eqref{eq:ex-lattice-clock-cap-receipt}:
\begin{equation}
 \frac{\zeta(j)}{\widehat\beta_j}\longrightarrow1,\qquad
 \frac{\zeta(j)}{\zeta(\operatorname{summit}(j))}\longrightarrow1
 \quad(j\to\infty),\qquad
 0\le\log\mathsf b_i\le(\log\widehat\beta_{j_i})^{48}
 \quad\text{eventually}.
 \label{eq:ex-lattice-scribe-clock-inputs}
\end{equation}
The definition of the rate gives
\[
 \frac{\log(r_i^{-1})}{\widehat\beta_{j_i}}
 =\frac{\log(\mathsf a_{j_i}^{-1})}{\widehat\beta_{j_i}}
   -\frac{\log\mathsf b_i}{\widehat\beta_{j_i}}
 =\frac1\kappa\left(1-\frac{\bar\beta}{\widehat\beta_{j_i}}\right)
   -\frac{\log\mathsf b_i}{\widehat\beta_{j_i}}
 \longrightarrow\frac1\kappa,
\]
because $(\log t)^{48}/t\to0$ as $t\to\infty$. Dividing by
$\zeta(j_i)/\widehat\beta_{j_i}\to1$ and multiplying by
$\zeta(j_i)/\zeta(k_i)\to1$ gives
\begin{equation}
 \frac{\log(r_i^{-1})}{q_i}\longrightarrow\frac1\kappa.
 \label{eq:ex-lattice-scribe-clock}
\end{equation}
In particular the slope cancels in the logarithmically normalized
trace. For the selected phase let
$Z_i=r_i^4U_{i,4}^{-1}\Theta$ be its pulled-back centered trace. The
physical transfer gives $q_iZ_i\to4d_2\mathcal O=-\kappa\mathcal O$,
whereas \eqref{eq:ex-lattice-scribe-clock} gives
$\log(r_i^{-1})/q_i\to1/\kappa$. Therefore
\begin{equation}\label{eq:scribe-slope-cancellation}
       \log(r_i^{-1})Z_i\longrightarrow-\mathcal O.
\end{equation}
At two separated points, the moving-label correlation comparison and
\eqref{eq:trace-reference-normalization} consequently give the
coefficient $S_0(-\mathcal O,-\mathcal O)=3\dim\mathfrak g/\pi^4$
at unit separation. This is the coefficient in the trace asymptotic;
using the same positive slope in the two factors is sufficient.

Put $\ell_i=\log(r_i^{-1})$. Since $r_i\to0$, $\ell_i\to\infty$, and
$\ell_ir_i^4=\ell_ie^{-4\ell_i}\to0$. For large $i$, $\ell_i>0$ and
\[
 q_ir_i^4=\frac{\ell_ir_i^4}{\ell_i/q_i}\longrightarrow\frac0{1/\kappa}=0,
\]
which is \eqref{eq:ex-lattice-scribe-band-cost}. Finally
$q_i\to\infty$, since $q_i=\ell_i/(\ell_i/q_i)$ and $\ell_i\to\infty$.

More explicitly, the scalar representation is
\[
 \iota(e,d)=e\,1+4d\,\mathcal O.
\]
The finite reference calculation shows that $1$ and $\mathcal O$ are
linearly independent, so $\iota$ is injective and has a continuous left
inverse on its range.
The limiting comparison tensor is $4d_2\mathcal O$, where
$d_2=-\kappa/4$, so its coordinates are $(0,d_2)$. Put
$Z_i=q_i r_i^4U_i^{-1}\Theta_T$. For the coordinate representation,
and only for it, assume in addition that $T$ carries a hierarchy and
respects $S$, and that the selected algebra has the local Ward structure
(conservation, symmetry and reality of its energy tensor and its tensor
transformation law). Then
$Z_i\in\operatorname{range}(\iota)$ eventually, separately from the limit
$Z_i\to4d_2\mathcal O$. Indeed, under these assumptions $\Theta_T$ is
fixed by the real structure of $S$ and by every signed permutation of the
axes. On the tail where the grade-four chart is coherent, $U_i^{-1}$
intertwines the real structures and the signed-permutation actions of
the two tables, and multiplication by the positive real number
$q_ir_i^4$ preserves both conditions; so $Z_i$ is a real,
signed-invariant grade-four field of the reference table. The diagonal
tensor with all four diagonal entries equal to $Z_i$ and vanishing
off-diagonal entries is then a real, symmetric, signed-covariant stress
tensor of the reference table, hence a real combination of the four basis
tensors of Lemma~\ref{lem:ex-trace-free-stress-basis}. Taking its trace,
and using that the trace of each of those basis tensors lies in
$\mathbb R1+\mathbb R\mathcal O$, shows
$4Z_i\in\mathbb R1+\mathbb R\mathcal O$. On that tail let $(a_i,b_i)=\iota^{-1}(Z_i)$;
continuity of the fixed left inverse gives $a_i\to0$ and $b_i\to d_2$.
Define $e_i=a_i/q_i$ and $d_i=b_i/q_i$. Thus, eventually,
\[
 r_i^4U_i^{-1}\Theta_T=\iota(e_i,d_i),
 \qquad q_i e_i\to0,\qquad q_i d_i\to d_2.
\]
The definition of the zoom then gives
$q_i r_i^4U_i^{-1}\Theta_T\to4d_2\mathcal O$.

The exact record contains $q_i\to\infty$ and the clock limit
\eqref{eq:ex-lattice-scribe-clock}, in which the scale is the reciprocal
$r_i^{-1}$ of the chart radius and not the zoom factor;
the same finite-volume
chart and hierarchy; the physical stress limit; for every table $T$ the
zoomed trace limit; and, for every table $T$ that carries a hierarchy and
respects $S$, the late scalar representation and two coefficient limits
just proved. The fixed-scale case uses the selected
chart directly. For arbitrary positive null sequences the order of choices
is as follows. The ultrafilter, the arch scale and the reference hierarchy
are fixed first. Then, for every positive sequence $s_n\to0$, the paid
scale selection chooses a strictly increasing sequence of indices $n_i$
and a scaling whose radii are exactly $r_i=s_{n_i}$
(Lemma~\ref{lem:ex-lattice-firesky-selection}, using the joint
chart purchase proved in that section); the same history is then passed to the
finite-volume chart producer, and the physical transfer above applies
without another extraction. With the actual shell, Ward comparison-bound and branch-recording instances
specified above, applying the response, Ward comparison, branch and scale-selection
producers to this same history gives the complete trace chart data. The Ward comparison's marked-vector limits, the actual
arch identification and clock limits, and the physical-energy
normalization have been proved above on that history. The
tested-bound, adjacent-column and actual branch-recording arguments
supply the additional quantitative instances for the external Ward comparison
and branch theorems. The
compact comparison and the shell limit are consequences of the
selected chart family's uniform short-test bound and the derivative laws,
as proved above; they are not additional independent inputs.
